\section{Exact and approximate gluing}\label{sec:gluing}

\subsection{The sheaf used by the development}
Let $X$ be a topological space with the measurable and Borel structure required
by the formal statement. Let $F_0$ assign to an open subset its finite spatial
measures, with restriction maps. The Lean definition
\texttt{probabilityPresheaf} denotes the sheafification $F$ of $F_0$.
Despite that identifier, the initial local data are unnormalised finite
measures; mass one is not part of this definition.

\begin{claim}
$F$ is a sheaf.
\end{claim}
\begin{proof}
Sheafification returns a sheaf together with its sheaf-property proof.
Projecting that proof establishes the claim for $F$.
\end{proof}
\formal{sheaf-property}
\scope This result uses the sheafification construction. It supplies no
physical mechanism for producing compatible local data.

\subsection{Compatible local sections have a unique extension}
\begin{claim}
Let $F$ be any sheaf of types on $X$, and let $(U_i)_{i\in I}$ be an open
cover of $X$. Suppose $s_i\in F(U_i)$ and
\[
 s_i|_{U_i\cap U_j}=s_j|_{U_i\cap U_j}\qquad(i,j\in I).
\]
Then exactly one $g\in F(X)$ satisfies $g|_{U_i}=s_i$ for every $i$.
\end{claim}
\begin{proof}
The unique-gluing form of the sheaf condition gives a section on
$\bigcup_i U_i$, with the specified restrictions, unique among sections there.
The cover equality identifies that union with $X$. Transport the section along
this equality. Functoriality of the restriction maps shows that its restriction
to each $U_i$ is $s_i$.

Conversely, transport any competing section on $X$ back to the union. Its
restrictions are again the $s_i$, so the uniqueness supplied by the sheaf
condition identifies it with the first section. Transporting forward gives the
required equality on $X$. In the Lean proof, equality of the relevant inclusion
maps uses the fact that morphisms between opens form a subsingleton.
\end{proof}
\formal{sheaf-gluing}

\begin{claim}
The same conclusion holds for the development's sheafified spatial measures.
\end{claim}
\begin{proof}
Apply the preceding claim to $F=\texttt{probabilityPresheaf}$ and use the
sheaf property established above. The compatible family and cover equality
are exactly the hypotheses of this application.
\end{proof}
\formal{probability-gluing}
\scope This proof is a direct term applying a previous theorem. It has no
tactic history; the extracted statement and proof dependencies still cover it.

\subsection{What equilibrium contributes}
A thermodynamic cover supplies local sections $s_i$, phases $\theta_i$, and a
finite, symmetric, strictly positive coupling matrix. Its structure includes
two substantive conditions: the phase configuration globally minimises the
reduced Kuramoto potential, and
\[
 \cos(\theta_i-\theta_j)=1
 \quad\Longrightarrow\quad
 s_i|_{U_i\cap U_j}=s_j|_{U_i\cap U_j}.
\]
The second condition is an obligation on the supplied local data.

\begin{claim}
A thermodynamic cover determines a unique global section with its prescribed
local restrictions.
\end{claim}
\begin{proof}
The theorem characterising global minima of the reduced Kuramoto potential
implies phase locking for the cover: $\cos(\theta_i-\theta_j)=1$ for every pair.
The overlap-agreement field then makes the local sections compatible. Apply
the spatial-measure gluing theorem to the cover and that compatible family.
\end{proof}
\formal{thermodynamic-gluing}
\scope This argument consumes the supplied equilibrium and overlap conditions.
The existence of a physical system satisfying them is a separate matter.

\subsection{Approximate agreement gives a controlled selection}
Let the site set and patch index set be finite. Let $s_i(x)\geq0$ be local
mass profiles represented as functions on the ambient sites. Let
$w_i(x)\geq0$ satisfy $w_i(x)=0$ outside $U_i$ and
$\sum_i w_i(x)=1$ at every site. Assume that for $x\in U_i\cap U_j$,
$|s_i(x)-s_j(x)|\leq\varepsilon$, where $\varepsilon\geq0$.

\begin{claim}
The selected profile $g(x)=\sum_j w_j(x)s_j(x)$ satisfies
$\sup_{x\in U_i}|g(x)-s_i(x)|\leq\varepsilon$ for every patch $i$.
\end{claim}
\begin{proof}
Fix $x\in U_i$. Every nonzero weight $w_j(x)$ has $x\in U_j$, so its profile
is within $\varepsilon$ of $s_i(x)$. Normalisation and the triangle inequality
therefore give
\[
 |g(x)-s_i(x)|
 =\left|\sum_j w_j(x)(s_j(x)-s_i(x))\right|
 \leq\sum_j w_j(x)\varepsilon=\varepsilon.
\]
The terms with zero weight contribute zero. The uniform metric on the finite
patch converts these pointwise bounds into the displayed profile-distance
bound. Lean isolates the weighted inequality in its helper
\texttt{weighted\_dist\_le} and then applies it at each site.
\end{proof}
\formal{approximate-gluing}
\scope The profiles need not have total mass one. The result provides an
approximation in the specified uniform mass metric, and assumes the overlap
bounds and the subordinate normalised weights.
